\documentclass[aspectratio=169]{beamer}

\usetheme{Madrid}
\usecolortheme{default}
\usefonttheme{professionalfonts}

\usepackage[T1]{fontenc}
\usepackage{lmodern}
\usepackage{amsmath,amssymb}
\usepackage{bm}
\usepackage{xcolor}

\title[PQC Hard Problem Colloquium]{Orientation Materials for the PQC Hard Problem Colloquium}
\author{Your Name}
\institute{Your Institution / Research Group}
\date{\today}

\begin{document}
	
	\begin{frame}
		\titlepage
	\end{frame}
	
	\begin{frame}{Purpose of the Colloquium}
		\begin{itemize}
			\item This colloquium studies the computational assumptions underlying classical and post-quantum cryptography.
			\item Our focus is not merely on cryptosystems, but on the \textbf{hard problems} that sustain them.
			\item For each family, we ask three questions:
			\begin{enumerate}
				\item What is the underlying hard problem?
				\item What are the best known attacks?
				\item How do classical and quantum complexities compare?
			\end{enumerate}
		\end{itemize}
	\end{frame}
	
	\begin{frame}{Why This Topic Matters}
		\begin{block}{Central Observation}
			Cryptography survives only as long as its underlying hardness assumptions survive.
		\end{block}
		
		\begin{itemize}
			\item Classical public-key cryptography relied mainly on factorization and discrete logarithms.
			\item Quantum algorithms, especially Shor's algorithm, fundamentally change this landscape.
			\item Post-quantum cryptography therefore asks which alternative problems remain hard even against quantum adversaries.
		\end{itemize}
	\end{frame}
	
	\section{Classical Foundations and Their Quantum Collapse}
	
	\begin{frame}{Classical Number-Theoretic Foundations}
		Historically, public-key cryptography was built on three dominant assumptions:
		\begin{itemize}
			\item \textbf{Integer Factorization} \hfill RSA, Rabin
			\item \textbf{Discrete Logarithm in Finite Fields} \hfill Diffie--Hellman, DSA, ElGamal
			\item \textbf{Elliptic Curve Discrete Logarithm} \hfill ECDH, ECDSA, EdDSA
		\end{itemize}
		
		\vspace{0.5em}
		These problems were attractive because they supported efficient and elegant public-key constructions.
	\end{frame}
	
	\begin{frame}{Integer Factorization}
		\begin{itemize}
			\item Integer factorization underlies \textbf{RSA} and \textbf{Rabin}.
			\item Classically, the best known attacks include:
			\begin{itemize}
				\item \textbf{Quadratic Sieve},
				\item \textbf{General Number Field Sieve}.
			\end{itemize}
			\item These are \textbf{sub-exponential} classical attacks.
			\item In the quantum setting, \textbf{Shor's algorithm} solves the problem in polynomial time.
		\end{itemize}
		
		\begin{block}{Orientation Point}
			Factorization is a model example of a problem that is hard enough classically, yet fundamentally broken quantumly.
		\end{block}
	\end{frame}
	
	\begin{frame}{Discrete Logarithm and Elliptic Curve DLP}
		\textbf{Finite-field DLP}
		\begin{itemize}
			\item Supports Diffie--Hellman, DSA, and ElGamal.
			\item Best classical attacks include \textbf{Index Calculus} and \textbf{Number Field Sieve}.
			\item These are again sub-exponential classically.
		\end{itemize}
		
		\vspace{0.5em}
		\textbf{Elliptic-curve DLP}
		\begin{itemize}
			\item Supports ECDH, ECDSA, and EdDSA.
			\item Generic classical attacks are roughly \textbf{$O(\sqrt{n})$}, e.g.\ Pollard's rho.
			\item Yet both finite-field DLP and ECDLP fall to \textbf{Shor's algorithm} in polynomial time.
		\end{itemize}
	\end{frame}
	
	\begin{frame}{The Transition to PQC}
		\begin{block}{Key Transition}
			The post-quantum transition is the replacement of Shor-vulnerable assumptions with alternative assumptions for which no efficient quantum algorithms are known.
		\end{block}
		
		\begin{itemize}
			\item This is not merely a standards transition.
			\item It is a transition in \textbf{mathematical foundations}.
			\item The main candidate families are:
			\begin{itemize}
				\item lattices,
				\item codes,
				\item hash-based constructions,
				\item multivariate systems,
				\item isogenies.
			\end{itemize}
		\end{itemize}
	\end{frame}
	
	\section{Main Post-Quantum Hard Problem Families}
	
	\begin{frame}{Lattice Problems}
		The lattice family includes:
		\begin{itemize}
			\item \textbf{LWE},
			\item \textbf{SIS},
			\item \textbf{SVP},
			\item \textbf{CVP}.
		\end{itemize}
		
		Representative systems include:
		\begin{itemize}
			\item Kyber / ML-KEM,
			\item Dilithium / ML-DSA,
			\item NTRU,
			\item fully homomorphic encryption constructions.
		\end{itemize}
	\end{frame}
	
	\begin{frame}{How Lattice Systems Are Attacked}
		The dominant attack language for lattice-based cryptography includes:
		\begin{itemize}
			\item \textbf{Lattice reduction}, especially LLL and BKZ,
			\item \textbf{Sieve algorithms},
			\item careful \textbf{concrete security estimation}.
		\end{itemize}
		
		\vspace{0.5em}
		The best known attacks remain broadly \textbf{exponential} in the relevant dimension parameter.
		
		\begin{block}{Orientation Point}
			Lattice cryptography is currently the most important practical PQC family precisely because no Shor-style collapse is known.
		\end{block}
	\end{frame}
	
	\begin{frame}{Code-Based Problems}
		Code-based cryptography is centered on:
		\begin{itemize}
			\item \textbf{syndrome decoding},
			\item more generally, decoding random linear codes.
		\end{itemize}
		
		Representative systems include:
		\begin{itemize}
			\item McEliece,
			\item HQC,
			\item BIKE.
		\end{itemize}
		
		The dominant attack family is \textbf{Information Set Decoding (ISD)} and its variants.
	\end{frame}
	
	\begin{frame}{Code-Based Security Perspective}
		\begin{itemize}
			\item Best known attacks remain \textbf{exponential}.
			\item Quantum improvements are comparatively limited.
			\item Grover's algorithm gives only a \textbf{quadratic speedup}, rather than a full structural break.
		\end{itemize}
		
		\begin{block}{Orientation Point}
			Code-based cryptography remains one of the most conservative and durable post-quantum directions.
		\end{block}
	\end{frame}
	
	\begin{frame}{Hash-Based Assumptions}
		Hash-based cryptography relies on:
		\begin{itemize}
			\item collision resistance,
			\item preimage resistance,
			\item second-preimage resistance.
		\end{itemize}
		
		Representative signature systems include:
		\begin{itemize}
			\item SPHINCS+,
			\item XMSS,
			\item LMS.
		\end{itemize}
	\end{frame}
	
	\begin{frame}{Hash-Based Security Perspective}
		\begin{itemize}
			\item Collision attacks follow the \textbf{birthday bound}.
			\item Quantum preimage search is accelerated by \textbf{Grover's algorithm}.
			\item But there is no analogue of Shor's algorithm causing polynomial-time collapse.
		\end{itemize}
		
		\begin{block}{Orientation Point}
			Hash-based cryptography is often viewed as assumption-conservative, even when it is less attractive in size or efficiency.
		\end{block}
	\end{frame}
	
	\begin{frame}{Multivariate Problems}
		Multivariate cryptography is based on the difficulty of solving systems of quadratic equations over finite fields.
		
		\vspace{0.5em}
		This is commonly framed as the \textbf{MQ problem}.
		
		\vspace{0.5em}
		Representative systems include:
		\begin{itemize}
			\item UOV,
			\item MAYO,
			\item historically, Rainbow.
		\end{itemize}
	\end{frame}
	
	\begin{frame}{Multivariate Attack Perspective}
		Relevant attack language includes:
		\begin{itemize}
			\item \textbf{MinRank},
			\item \textbf{Kipnis--Shamir},
			\item \textbf{XL / Gr\"obner basis} methods,
			\item and scheme-specific structural attacks.
		\end{itemize}
		
		\vspace{0.5em}
		While the general problem remains hard-looking, particular constructions may fail because they reveal exploitable algebraic structure.
	\end{frame}
	
	\begin{frame}{The Rainbow Lesson}
		\begin{itemize}
			\item Rainbow was a prominent multivariate signature construction.
			\item It was broken by a polynomial-time structural attack.
			\item This teaches an important lesson:
		\end{itemize}
		
		\begin{block}{Critical Distinction}
			Family-level hardness does not automatically imply scheme-level security.
		\end{block}
	\end{frame}
	
	\begin{frame}{Isogeny Problems}
		Isogeny-based cryptography is among the most mathematically elegant areas in PQC.
		
		\vspace{0.5em}
		Representative names include:
		\begin{itemize}
			\item SQISign,
			\item CSIDH,
			\item historically, SIDH / SIKE.
		\end{itemize}
	\end{frame}
	
	\begin{frame}{The SIKE Lesson}
		\begin{itemize}
			\item SIKE was broken by the \textbf{Castryck--Decru attack}.
			\item This attack was polynomial time and completely changed the status of SIDH/SIKE.
			\item Other attack frameworks include meet-in-the-middle methods and quantum algorithms such as Kuperberg's algorithm in certain commutative settings.
		\end{itemize}
		
		\begin{block}{Orientation Point}
			Isogeny cryptography remains mathematically rich, but it must now be studied with exceptional caution regarding hidden structure.
		\end{block}
	\end{frame}
	
	\section{How to Read the Hard-Problem Landscape}
	
	\begin{frame}{Comparative Reading of the Landscape}
		The families divide naturally into four security narratives:
		\begin{itemize}
			\item \textbf{Factorization / DLP / ECDLP}: classically strong, quantumly broken.
			\item \textbf{Lattice / Code-based}: currently the strongest practical PQC foundations.
			\item \textbf{Hash-based}: conservative and structurally simple.
			\item \textbf{Multivariate / Isogeny}: mathematically compelling, but more exposed to structural fragility.
		\end{itemize}
	\end{frame}
	
	\begin{frame}{Three Questions for Every Scheme}
		For every construction discussed in this colloquium, we should ask:
		\begin{enumerate}
			\item What exact hard problem is assumed?
			\item What are the best known classical and quantum attacks?
			\item Does the scheme introduce exploitable structure beyond the underlying hard problem?
		\end{enumerate}
		
		\vspace{0.75em}
		This provides a disciplined framework for comparing PQC proposals.
	\end{frame}
	
	\begin{frame}{What the Hard-Problem Summary Really Teaches}
		The summary should not be read merely as a catalog of attacks.
		
		\vspace{0.5em}
		It should be read as a set of research principles:
		\begin{itemize}
			\item distinguish \textbf{problem hardness} from \textbf{scheme security},
			\item distinguish \textbf{generic attacks} from \textbf{structural attacks},
			\item compare \textbf{classical} and \textbf{quantum} attack models,
			\item and study broken designs as lessons in cryptographic methodology.
		\end{itemize}
	\end{frame}
	
	\section{Orientation for the Colloquium}
	
	\begin{frame}{Suggested Colloquium Flow}
		A natural progression for the colloquium is:
		\begin{enumerate}
			\item Classical number-theoretic assumptions and Shor's algorithm
			\item Lattice problems and reduction-based security
			\item Code-based assumptions and ISD
			\item Hash-based signatures as a conservative baseline
			\item Multivariate constructions and structural failure
			\item Isogeny cryptography after SIKE
			\item Cross-family comparison of security assumptions
		\end{enumerate}
	\end{frame}
	
	\begin{frame}{Expected Outcomes}
		By the end of the colloquium, participants should be able to:
		\begin{itemize}
			\item explain the main hard-problem families in PQC,
			\item describe the best known attacks against each family,
			\item understand why some assumptions remain credible,
			\item and evaluate a cryptosystem at the level of its actual security assumption.
		\end{itemize}
	\end{frame}
	
	\begin{frame}{Closing Perspective}
		\begin{block}{Final Perspective}
			Post-quantum cryptography is not simply the search for new algorithms. It is the search for mathematical problems whose hardness survives quantum computation, structural cryptanalysis, and concrete attack advances.
		\end{block}
		
		\vspace{0.75em}
		That is the central purpose of the PQC Hard Problem Colloquium.
	\end{frame}

	\section{Symbolic Formulation of the Hard Problems}
	
	\begin{frame}{A General Cryptographic Template}
		For each cryptographic family, we study a computational relation of the form
		\[
		\mathcal{I} \mapsto \mathcal{S},
		\]
		where:
		\begin{itemize}
			\item $\mathcal{I}$ is the public instance,
			\item $\mathcal{S}$ is the hidden witness, secret, or structure,
			\item and the adversary is asked to recover $\mathcal{S}$, or to solve a decisional variant.
		\end{itemize}
		
		\vspace{0.5em}
		Thus, a hard problem is typically one of the following:
		\[
		\text{Search: } \mathcal{A}(\mathcal{I}) = \mathcal{S},
		\qquad
		\text{Decision: } \mathcal{A}(\mathcal{I}) \in \{0,1\},
		\]
		with no efficient known algorithm for the relevant parameter regime.
	\end{frame}
	
	%------------------------------------------------
	
	\begin{frame}{Integer Factorization}
		\textbf{Instance.}
		Let
		\[
		N = pq,
		\]
		where $p,q$ are large primes. The public input is $N$.
		
		\vspace{0.4em}
		\textbf{Search problem.}
		Recover nontrivial factors:
		\[
		\text{Given } N, \text{ find } p,q \text{ such that } N=pq.
		\]
		
		\vspace{0.4em}
		\textbf{Equivalent viewpoint.}
		The hardness lies in inverting the map
		\[
		(p,q) \longmapsto pq.
		\]
		
		\vspace{0.4em}
		\textbf{Why cryptographically useful.}
		RSA relies on the asymmetry:
		\[
		\text{multiplication is easy}, \qquad \text{factorization is hard}.
		\]
		
		\vspace{0.4em}
		\textbf{Attack language.}
		\[
		\text{QS}, \quad \text{GNFS}, \quad \text{Shor}.
		\]
		
		\[
		L_N\!\left[\frac13,c\right]
		\quad \text{classically}, 
		\qquad
		\mathrm{poly}(\log N)
		\quad \text{quantumly}.
		\]
	\end{frame}
	
	%------------------------------------------------
	
	\begin{frame}{Discrete Logarithm Problem in Finite Fields}
		Let $G = \langle g \rangle$ be a cyclic group, typically
		\[
		G = \mathbb{F}_p^\times.
		\]
		
		\textbf{Instance.}
		Given
		\[
		g \in G, \qquad h = g^x \in G.
		\]
		
		\textbf{Search problem.}
		Recover the exponent $x$:
		\[
		\text{find } x \text{ such that } h=g^x.
		\]
		
		\vspace{0.4em}
		\textbf{Asymmetry.}
		Exponentiation is easy:
		\[
		x \mapsto g^x,
		\]
		but inversion is believed hard.
		
		\vspace{0.4em}
		\textbf{Diffie--Hellman relation.}
		Given
		\[
		g^a,\, g^b,
		\]
		compute
		\[
		g^{ab}.
		\]
		
		\vspace{0.4em}
		\textbf{Attack language.}
		\[
		\text{Index Calculus}, \quad \text{NFS}, \quad \text{Shor}.
		\]
		
		\[
		L_p\!\left[\frac13,c\right]
		\quad \text{classically}, 
		\qquad
		\mathrm{poly}(\log p)
		\quad \text{quantumly}.
		\]
	\end{frame}
	
	%------------------------------------------------
	
	\begin{frame}{Elliptic Curve Discrete Logarithm Problem (ECDLP)}
		Let $E/\mathbb{F}_q$ be an elliptic curve, and let
		\[
		P \in E(\mathbb{F}_q)
		\]
		be a point of large order $n$.
		
		\textbf{Instance.}
		Given
		\[
		P, \qquad Q = xP.
		\]
		
		\textbf{Search problem.}
		Recover
		\[
		x \in \mathbb{Z}_n
		\quad \text{such that} \quad
		Q = xP.
		\]
		
		\vspace{0.4em}
		\textbf{Group law formulation.}
		The map
		\[
		x \longmapsto xP
		\]
		is easy to evaluate, but hard to invert.
		
		\vspace{0.4em}
		\textbf{Generic attack baseline.}
		Pollard's rho yields roughly
		\[
		O(\sqrt{n}).
		\]
		
		\vspace{0.4em}
		\textbf{Exceptional structure.}
		For special curves, one may reduce ECDLP to finite-field DLP via pairings:
		\[
		E(\mathbb{F}_q) \rightsquigarrow \mathbb{F}_{q^k}^{\times}.
		\]
		
		\vspace{0.4em}
		\textbf{Quantumly.}
		Shor again gives polynomial-time recovery of $x$.
	\end{frame}
	
	%------------------------------------------------
	
	\begin{frame}{Lattice Problems: Geometric Form}
		A lattice generated by basis vectors $\bm{b}_1,\dots,\bm{b}_m \in \mathbb{R}^n$ is
		\[
		\Lambda = \left\{ \sum_{i=1}^m z_i \bm{b}_i : z_i \in \mathbb{Z} \right\}.
		\]
		
		\textbf{Shortest Vector Problem (SVP).}
		Find
		\[
		\bm{v} \in \Lambda \setminus \{\bm{0}\}
		\]
		minimizing
		\[
		\|\bm{v}\|.
		\]
		
		\textbf{Closest Vector Problem (CVP).}
		Given $\bm{t} \in \mathbb{R}^n$, find
		\[
		\bm{v} \in \Lambda
		\]
		minimizing
		\[
		\|\bm{t}-\bm{v}\|.
		\]
		
		\vspace{0.4em}
		\textbf{Interpretation.}
		These are geometric optimization problems over a discrete subgroup of $\mathbb{R}^n$.
		
		\vspace{0.4em}
		\textbf{Attack language.}
		\[
		\text{LLL}, \quad \text{BKZ}, \quad \text{sieving}.
		\]
	\end{frame}
	
	%------------------------------------------------
	
	\begin{frame}{Learning With Errors (LWE)}
		LWE is an average-case algebraic-noise problem.
		
		\textbf{Secret.}
		\[
		\bm{s} \in \mathbb{Z}_q^n.
		\]
		
		\textbf{Samples.}
		For random
		\[
		\bm{a}_i \in \mathbb{Z}_q^n,
		\]
		we observe
		\[
		b_i = \langle \bm{a}_i,\bm{s} \rangle + e_i \pmod q,
		\]
		where $e_i$ is a small error.
		
		\vspace{0.4em}
		\textbf{Search-LWE.}
		Given many pairs
		\[
		(\bm{a}_i,b_i),
		\]
		recover $\bm{s}$.
		
		\textbf{Decision-LWE.}
		Distinguish
		\[
		(\bm{a}_i,\langle \bm{a}_i,\bm{s}\rangle + e_i)
		\]
		from uniform samples in
		\[
		\mathbb{Z}_q^n \times \mathbb{Z}_q.
		\]
		
		\vspace{0.4em}
		\textbf{Symbolic intuition.}
		Without noise, one solves a linear system.
		With noise, the system becomes
		\[
		A\bm{s} + \bm{e} \equiv \bm{b} \pmod q,
		\]
		and the small perturbation $\bm{e}$ is exactly what creates hardness.
	\end{frame}
	
	%------------------------------------------------
	
	\begin{frame}{SIS: Short Integer Solution}
		Let
		\[
		A \in \mathbb{Z}_q^{m \times n}.
		\]
		
		\textbf{Problem.}
		Find a nonzero vector
		\[
		\bm{x} \in \mathbb{Z}^n
		\]
		such that
		\[
		A\bm{x} \equiv \bm{0} \pmod q
		\]
		and
		\[
		\|\bm{x}\|
		\]
		is small.
		
		\vspace{0.4em}
		\textbf{Interpretation.}
		SIS asks for a short relation among the columns of $A$ modulo $q$.
		
		\vspace{0.4em}
		\textbf{Why important.}
		SIS is a canonical source of hash functions and signature constructions.
		
		\vspace{0.4em}
		\textbf{Symbolic contrast with LWE.}
		\[
		\text{LWE: hide } \bm{s} \text{ with noise},
		\qquad
		\text{SIS: find short dependence modulo } q.
		\]
	\end{frame}
	
	%------------------------------------------------
	
	\begin{frame}{Code-Based Hardness: Syndrome Decoding}
		Let
		\[
		H \in \mathbb{F}_2^{r \times n}
		\]
		be a parity-check matrix.
		
		\textbf{Given.}
		A syndrome
		\[
		\bm{s} \in \mathbb{F}_2^r
		\]
		and a weight bound $w$.
		
		\textbf{Problem.}
		Find an error vector
		\[
		\bm{e} \in \mathbb{F}_2^n
		\]
		such that
		\[
		H\bm{e}^T = \bm{s}^T
		\]
		and
		\[
		\mathrm{wt}(\bm{e}) \le w.
		\]
		
		\vspace{0.4em}
		\textbf{Interpretation.}
		The vector $\bm{e}$ is sparse, but hidden by the linear map $H$.
		
		\vspace{0.4em}
		\textbf{Cryptographic asymmetry.}
		Computing
		\[
		\bm{s} = H\bm{e}^T
		\]
		is easy; recovering a low-weight $\bm{e}$ is hard.
		
		\vspace{0.4em}
		\textbf{Attack language.}
		\[
		\text{ISD},\ \text{Stern},\ \text{Lee--Brickell},\ \text{BJMM-type methods}.
		\]
	\end{frame}
	
	%------------------------------------------------
	
	\begin{frame}{Hash-Based Hardness}
		Let
		\[
		H : \{0,1\}^\ast \to \{0,1\}^n
		\]
		be a cryptographic hash function.
		
		\textbf{Preimage resistance.}
		Given
		\[
		y \in \{0,1\}^n,
		\]
		find
		\[
		x \text{ such that } H(x)=y.
		\]
		
		\textbf{Second-preimage resistance.}
		Given $x$, find
		\[
		x' \neq x
		\quad \text{such that} \quad
		H(x') = H(x).
		\]
		
		\textbf{Collision resistance.}
		Find
		\[
		x \neq x'
		\quad \text{such that} \quad
		H(x)=H(x').
		\]
		
		\vspace{0.4em}
		\textbf{Generic bounds.}
		\[
		\text{collision search: } O(2^{n/2}),
		\qquad
		\text{preimage search: } O(2^n).
		\]
		
		\textbf{Quantum effect.}
		Grover reduces exhaustive preimage search to roughly
		\[
		O(2^{n/2}).
		\]
	\end{frame}
	
	%------------------------------------------------
	
	\begin{frame}{Multivariate Quadratic (MQ) Problem}
		Let
		\[
		f_1,\dots,f_m \in \mathbb{F}_q[x_1,\dots,x_n]
		\]
		be multivariate quadratic polynomials.
		
		\textbf{Problem.}
		Solve the system
		\[
		f_1(x_1,\dots,x_n)=\cdots=f_m(x_1,\dots,x_n)=0.
		\]
		
		\vspace{0.4em}
		Equivalently, for
		\[
		\bm{x}=(x_1,\dots,x_n),
		\]
		find
		\[
		\bm{x} \in \mathbb{F}_q^n
		\]
		such that
		\[
		F(\bm{x})=\bm{0}.
		\]
		
		\vspace{0.4em}
		\textbf{Why difficult.}
		Quadratic interactions create nonlinear constraints:
		\[
		x_ix_j,\quad x_i^2.
		\]
		
		\vspace{0.4em}
		\textbf{Attack language.}
		\[
		\text{XL}, \quad \text{F4/F5}, \quad \text{Gr\"obner bases}, \quad \text{MinRank reductions}.
		\]
		
		\vspace{0.4em}
		\textbf{Important warning.}
		General MQ hardness and the security of a particular MQ scheme are not the same statement.
	\end{frame}
	
	%------------------------------------------------
	
	\begin{frame}{Isogeny Problems}
		Let
		\[
		E_1, E_2
		\]
		be elliptic curves over a finite field.
		
		\textbf{Basic object.}
		An isogeny is a nonconstant morphism
		\[
		\phi : E_1 \to E_2
		\]
		that is also a group homomorphism.
		
		\vspace{0.4em}
		\textbf{Hardness intuition.}
		Given information about the endpoints, recover hidden isogeny data:
		\[
		E_1 \rightsquigarrow E_2.
		\]
		
		Depending on the system, the task may be:
		\begin{itemize}
			\item finding a path in an isogeny graph,
			\item recovering an ideal-class action,
			\item or reconstructing auxiliary hidden structure.
		\end{itemize}
		
		\vspace{0.4em}
		\textbf{Symbolic viewpoint.}
		The public data describes vertices in a graph,
		\[
		\mathcal{G}_{\mathrm{iso}},
		\]
		while the secret corresponds to a hidden path, action, or morphism.
		
		\vspace{0.4em}
		\textbf{Caution.}
		Some specific formulations, notably SIDH/SIKE, were broken by exploiting extra structure not captured by the naive hardness picture.
	\end{frame}
	
	%------------------------------------------------
	
	\begin{frame}{A Unifying View}
		Each family can be summarized as ``easy forward map, hard inverse problem'':
		\[
		\begin{aligned}
			(p,q) &\mapsto pq && \text{(Factorization)} \\
			x &\mapsto g^x && \text{(DLP)} \\
			x &\mapsto xP && \text{(ECDLP)} \\
			(\bm{s},\bm{e}) &\mapsto (A, A\bm{s}+\bm{e}) && \text{(LWE)} \\
			\bm{x} &\mapsto A\bm{x} \bmod q && \text{(SIS / modular relations)} \\
			\bm{e} &\mapsto H\bm{e}^T && \text{(Syndrome decoding)} \\
			x &\mapsto H(x) && \text{(Hash-based)} \\
			\bm{x} &\mapsto F(\bm{x}) && \text{(MQ)} \\
			\phi &\mapsto \text{public isogeny data} && \text{(Isogenies)}
		\end{aligned}
		\]
		
		\vspace{0.5em}
		The colloquium is therefore about one central theme:
		\[
		\textbf{Which inverse problems remain hard in the quantum era?}
		\]
	\end{frame}
	
	
	\begin{frame}
		\centering
		{\LARGE Thank you.}
		
		\vspace{0.5em}
		{\large Questions and discussion}
	\end{frame}
	
\end{document}